% Rubbish 投稿模板（LaTeX 版）· 2026-10
% 使用说明：xelatex 编译（中文需要）；替换「」内容后删除注释行。
\documentclass[12pt,a4paper]{article}
\usepackage[margin=2.5cm]{geometry}
\usepackage{ctex}          % 中文支持（xelatex）
\usepackage{graphicx}
\usepackage{booktabs}
\usepackage{amsmath}
\usepackage{hyperref}

\title{\textbf{「论文标题：一句话说清你的研究」}\\
{\large 正经赛道（Serious Track）或荒诞赛道（Absurd Track）——二选一删另一个}}
\author{「第一作者」\and 「合作者」}
\date{「投稿日期」}

\begin{document}
\maketitle

\begin{abstract}
「150--300 字。正经研究按『问题---方法---结果---结论』四段式；
荒诞研究同样要有完整论证闭环。」\\[1em]
\textbf{关键词：}「3--6 个，分号分隔」
\end{abstract}

\section{引言}
「研究背景与动机。正经研究交代学术脉络与缺口；荒诞研究交代现象的
火爆程度与『严肃对待它』的理由。」

\section{相关工作}
「综述已有文献或该现象的传播脉络，说明本文与其差异。」

\section{方法}
「数据来源、实验设计、分析方法，或荒诞研究的论证框架。
涉及人的调查请说明知情同意。」

\section{结果}
「核心发现。图表要求：自明、有编号与标题、正文中必须被引用。」

% 表格示例
\begin{table}[htbp]
\centering
\caption{「表标题」}
\begin{tabular}{lcc}
\toprule
「条件」 & 「指标 A」 & 「指标 B」 \\
\midrule
「对照组」 & 0.00 & 0.00 \\
「实验组」 & 0.00 & 0.00 \\
\bottomrule
\end{tabular}
\end{table}

% 公式示例（正经赛道常用；荒诞赛道用来说明『严谨地荒诞』）
\begin{equation}
R = 100 \times \left( 0.5\,r + 0.3\,d + 0.2\,s \right)
\end{equation}

\section{讨论}
「结果意味着什么？局限是什么？荒诞研究请诚实讨论『荒诞在哪里』。」

\section{结论}
「3--5 句收尾：结论、局限、后续工作。」

\section*{致谢}
「基金资助、AI 辅助声明（工具与用途）、提供帮助的人。」

\begin{thebibliography}{9}
\bibitem{ref1} 「作者. 题目. 刊物, 年份.」
\end{thebibliography}

\end{document}
